% Generated by scripts/assemble.py; corrections belong in source/corrections.json.
\examhead{105}{中国科学技术大学 2025 年数学分析试卷}{104}
% p105-b01; source PDF page 105
\par\Needspace{5\baselineskip}
\blockanchor{p105-b01}
\majorquestion{一}计算题. 每题 10 分, 共 40 分.\par
% p105-b02; source PDF page 105
\blockanchor{p105-b02}
(1) 计算 $\lim_{x\to \infty}\left(\cos \frac{1}{x}\right)^{x^2}.$\par
% p105-b03; source PDF page 105
\blockanchor{p105-b03}
(2) 计算 $\lim_{n\to \infty}\frac{\sqrt[n]{n!}}{n}$.\par
% p105-b04; source PDF page 105
\blockanchor{p105-b04}
(3) 计算 $\lim_{n\to \infty}\int_0^{+\infty}x^2 e^{-x^{2n}}\mathrm{d}x.$\par
% p105-b05; source PDF page 105
\blockanchor{p105-b05}
(4) 由方程 $\sin x + \ln y - xy^3 = 0$ 在 $(0,1)$ 附近确定的隐函数 $y = f(x)$, 求 $f'(0)$.\par
% p105-b06; source PDF page 105
\par\Needspace{5\baselineskip}
\par\Needspace{6\baselineskip}
\blockanchor{p105-b06}
\majorquestion{二}(15 分) 把 $f(x) = \frac{\pi}{2} - x, x \in [0, \pi]$ 展开成余弦级数, 并计算\par
\nopagebreak[4]
% p105-b07; source PDF page 105
\blockanchor{p105-b07}
\displayformula{\sum_ {n = 1} ^ {\infty} \frac {1}{(2 n - 1) ^ {2}}, \sum_ {n = 1} ^ {\infty} \frac {1}{n ^ {4}}, \sum_ {n = 1} ^ {\infty} \frac {1}{n (n + 1)} \left(1 + \frac {1}{2} + \dots + \frac {1}{n}\right).}
% p105-b08; source PDF page 105
\par\Needspace{5\baselineskip}
\par\Needspace{6\baselineskip}
\blockanchor{p105-b08}
\majorquestion{三}(15 分) $f$ 是定义在 $D = \{(x, y) \mid x^2 + y^2 \leq \pi\}$ 上的二阶连续可微函数, 且\par
\nopagebreak[4]
% p105-b09; source PDF page 105
\blockanchor{p105-b09}
\displayformula{\frac {\partial^ {2} f}{\partial x ^ {2}} + \frac {\partial^ {2} f}{\partial y ^ {2}} = \sin (x ^ {2} + y ^ {2}).}
% p105-b10; source PDF page 105
\blockanchor{p105-b10}
计算 $\oint_{\partial D} \frac{\partial f}{\partial n} \, \mathrm{d}s,$ 其中 $n$ 为 $\partial D$ 的单位外法向量.\par
% p105-b11; source PDF page 105
\par\Needspace{5\baselineskip}
\par\Needspace{6\baselineskip}
\blockanchor{p105-b11}
\majorquestion{四}(15 分) $\Sigma$ 为 $z = x^{2} + y^{2}$, $0 \leq z \leq 1$, 取下侧, 计算\par
\nopagebreak[4]
% p105-b12; source PDF page 105
\blockanchor{p105-b12}
\displayformula{\iint_ {\Sigma} (2 y + z) \mathrm{d} y \mathrm{d} z + z \mathrm{d} x \mathrm{d} y.}
% p105-b13; source PDF page 105
\par\Needspace{5\baselineskip}
\blockanchor{p105-b13}
\majorquestion{五}(15 分) 计算 $\int_{0}^{+\infty} e^{-\alpha x^2} \cos(\beta x) \, \mathrm{d}x$, 其中 $\alpha > 0, \beta \in \mathbb{R}$.\par
% p105-b14; source PDF page 105
\par\Needspace{5\baselineskip}
\blockanchor{p105-b14}
\majorquestion{六}(20 分) 定义在 $[0,1]$ 上的连续函数列 $\{f_n(x)\}_{n\in \mathbb{N}}$ 一致收敛于 $f(x)$.\par
% p105-b15; source PDF page 105
\blockanchor{p105-b15}
(1) 证明: 对任意的 $\varepsilon > 0$, 存在 $\delta > 0$, 使得当 $x, y \in [0,1]$ 且 $|x - y| < \delta$ 时, $|f_n(x) - f_n(y)| < \varepsilon$ 对 $n = 1,2,\cdots$ 成立.\par
% p105-b16; source PDF page 105
\blockanchor{p105-b16}
(2) 证明: 存在 $M > 0$, 使得 $\sup_{x \in [0,1]} \{|f_n(x)|\} \leq M$ 对所有 $n = 1,2,\cdots$ 成立.\par
% p105-b17; source PDF page 105
\par\Needspace{5\baselineskip}
\blockanchor{p105-b17}
\majorquestion{七}(20 分) 设 $f(x) = \sum_{n=2}^{\infty} \frac{x^n}{n^2 \ln n}$.\par
% p105-b18; source PDF page 105
\blockanchor{p105-b18}
(1) 证明: $f(x)$ 在 $[-1, 1]$ 上连续.\par
% p105-b19; source PDF page 105
\blockanchor{p105-b19}
(2) 证明: $f(x)$ 在 $x = -1$ 处右可导.\par
% p105-b20; source PDF page 105
\blockanchor{p105-b20}
(3) 证明: $\lim_{x \to 1^{-}} f'(x) = +\infty$.\par
% p105-b21; source PDF page 105
\blockanchor{p105-b21}
(4) 证明: $f(x)$ 在 $x = 1$ 处左不可导.\par
% p105-b22; source PDF page 105
\par\Needspace{5\baselineskip}
\blockanchor{p105-b22}
\majorquestion{八}(10 分) 证明: $\mathbb{R}$ 上任意函数的严格极大值点至多可数.\par
